\section{Discussion}
\label{sec:review-discussion}
\subsection{Economic interpretation}
\begin{mergedcontribution}{authorSowmya}
\subsubsection{O1: Concentration and ownership effects}
\textbf{From oligopoly towards asymmetric duopoly.} Ownership-HHI rose from a mean of 2,408 before the Jet Airways exit to 4,877 after the mergers, equivalent to falling from 4.2 to 2.1 equal-sized airlines; nine in ten passengers now fly IndiGo or the Air India Group. Concentration never left the highly concentrated band, so what changed is the distance above it, and it changed in steps tied to exits and changes of control rather than through gradual drift (Tables~\ref{tab:o1-periods}, \ref{tab:o1-drift} and~\ref{tab:o1-regression}).

\textbf{Ownership matters more than brands.} The Tata acquisition is the largest step (about 900 points), and half of it is the Grouping Effect: brands that remained separate in DGCA reports came under one owner. Brand shares alone would understate that step by half. Partial ownership is small: the wedge averages 23 points before 2022 and never exceeds 77 (Tables~\ref{tab:o1-gap} and~\ref{tab:review-o1-wedge}).

\textbf{Course lenses.} Read through structure--conduct--performance \parencite{gandhi2026marketstructure}, the rise signals greater potential market power. Under Cournot competition the industry Lerner index is $(\mathrm{HHI}/10{,}000)/|\varepsilon|$; with $|\varepsilon|=1.5$ the permanent drifts D1 and D5 imply margins 3.5 and 4.6 points higher (using the drift-point HHI changes in Table~\ref{tab:o1-drift}). This is an implication of structure, not a measured markup. IndiGo's large aircraft orders act as Stackelberg-style capacity commitments, and its share rose from 57.5\% to 63.2\% across D5 and was held. Both exits occurred with HHI already above 1,800, and Akasa Air's entry in August 2022 was followed only by a slow fall of about 390 points over six months, so entry has not offset consolidation.

\textbf{Supplementary grouping-gap comparison.} The gap is always $\geq0$. Its changes are $+316$ at D4 and $-100$ at D6, compared with $-27$ and $+29$ at the exits (Table~\ref{tab:o1-drift}). The wedge changes no month's concentration classification. The event steps remain significant with 6 or 18 lags and with COVID months dropped ($t\geq2.2$), and CR2 agrees (Tables~\ref{tab:o1-robust} and~\ref{tab:o1-cr2reg}). These are dated level shifts in national structure, not causal estimates of effects on fares, profits or passengers.
\end{mergedcontribution}

\begin{mergedcontribution}{authorMeena}
\subsubsection{O2: Exits and Tata consolidation}
\textbf{Exits.} Jet Airways' and Go First's exits reduced the effective number of equal-sized firms from 4.1 to 3.5 and 2.5 to 2.2, respectively. In the Cournot model, firms choose quantities simultaneously, taking rivals' output as given. Applied to our study, this would mean each airline decides how many seats to offer, assuming its competitors keep their seat numbers unchanged. Holding demand elasticity fixed, higher concentration implies a higher share-weighted markup:
\[
\overline{L} = \frac{\HHI/10{,}000}{|\varepsilon|},
\]
where $\overline{L}$ is the share-weighted price--cost markup as a fraction of price, and $|\varepsilon|$ is absolute market demand elasticity. An illustrative elasticity of 1.5 implies markup increases of about 3.0 and 4.6 percentage points of price, not observed fare increases: we have no fare data. The exit-window results are in Tables~\ref{tab:o2-absorption} and~\ref{tab:review-o2-capacity-contribution}. Worked calculations are in Appendix~\ref{app:workings}, part~C.

\textbf{Why IndiGo absorbed most.} In the Stackelberg model, a leader commits to its output first, anticipating how rivals will respond; followers then choose their output after observing that commitment. Applied to our study, this would mean IndiGo commits to future capacity first, and rivals plan their own expansion around that commitment. Its June 2023 order for 500 aircraft fits this interpretation, rather than proving the model describes actual behaviour. Slot access also constrains rivals' expansion: SpiceJet received vacated Jet Airways slots and expanded capacity, illustrating an entry barrier. IndiGo's dominance may reflect cost leadership as well as weaker competition; without fare and cost data, we cannot distinguish them.

\textbf{Tata consolidation and merger review.} For each set of entities combined, the fixed-share concentration effect is
\[
\Delta\HHI_{\text{mechanical}} = \left(\sum_i s_i\right)^2 - \sum_i s_i^2,
\]
where $s_i$ are pre-consolidation percentage shares, held fixed; effects are summed across separate mergers. Applied to our study, this means concentration can rise simply because airlines are counted as one entity rather than separate competitors, even if none gains any passengers. Common ownership accounts for most of the January 2022 rise; the late-2024 mergers mechanically added 256 brand \HHI{} points (worked calculations in Appendix~\ref{app:workings}, part~B). CCI's capacity commitments protected seven routes but did not address national concentration. Both perspectives matter: with two groups carrying about 90\% of passengers, national structure also affects slot access and exposure to single-airline disruption.
\end{mergedcontribution}

\begin{mergedcontribution}{authorSweta}
\subsubsection{O3: Entrant scale and financial viability}
\label{sec:review-o3-interpretation}
Nine entrants entered India's domestic aviation market over the eleven-and-a-half-year study period, but only two achieved meaningful and persistent national market shares. Vistara subsequently merged with Air India in November 2024, leaving Akasa Air as the only scale entrant continuing to expand as an independent competitor at the July 2026 cutoff. This suggests that entry alone did not ensure effective competition; entrants needed to build sufficient traffic and sustain operations over time (Tables~\ref{tab:o3-summary} and~\ref{tab:o3-routes}).

The findings are broadly consistent with the cost structure described by $AC=F/Q+v$. Higher passenger volumes can spread fixed costs over a larger output base, while better aircraft utilisation can improve the economics of operations. The load-factor evidence illustrates this difference: Vistara's load factor increased from 65.6\% in its first year to 84.7\% in its third, while Akasa Air recorded 84.7\% and 91.0\%, respectively. TruJet, by comparison, recorded 80.1\% in its first year and 76.6\% in its third. Although these figures are consistent with the potential benefits of higher utilisation, they do not establish that load factor or scale caused the different outcomes. TruJet's experience also shows why local route dominance is not equivalent to national scale: it held a dominant position on several routes, but the routes themselves carried relatively few passengers. A high share of a small market may therefore provide insufficient traffic to support the broader cost structure of an airline.

The financial results further qualify the relationship between growth and viability. Vistara's margins improved as revenue expanded, whereas Akasa Air's faster revenue growth was accompanied by continued substantial losses. Its continued expansion alongside these losses indicates the importance of access to external capital, but funding availability should not be interpreted as evidence of financial self-sufficiency. More generally, accounting results do not capture the full economic cost of capital; economic profitability would also require a return sufficient to compensate investors for the opportunity cost of their funds (Tables~\ref{tab:o3-fin} and~\ref{tab:o3-seven}).

Overall, the evidence supports the expectation that only entrants achieving sufficient scale and sustained passenger volumes became meaningful national competitors. It also shows that local route dominance and revenue growth alone were insufficient indicators of financial viability or survival. However, the analysis is descriptive: it identifies patterns consistent with the economic framework and existing literature, rather than testing whether economies of scale, funding or route concentration caused individual entrant outcomes.
\end{mergedcontribution}

\begin{mergedcontribution}{authorGoutham}
\subsubsection{O4: Concentration and resilience}
\textbf{National market structure.} During the focus period (December 2024 to July 2026), the national results are consistent with an asymmetric duopoly. IndiGo and Air India Group dominate the market, but IndiGo has a much larger passenger share (Tables~\ref{tab:o4-national-results} and~\ref{tab:o4-route-focus}). Other airlines form a small competitive fringe. This indicates that domestic passenger traffic depends heavily on two ownership groups, particularly IndiGo, despite the presence of several smaller airlines.

\textbf{Competition on individual airport pairs.} Table~\ref{tab:o4-route-focus} shows that IndiGo and Air India Group, the two national leaders, also dominate scheduled flights on the top 20 and top 10 busiest airport pairs. However, the median gap between the two largest groups is smaller than nationally. This highlights the importance of market definition: IndiGo's large national lead does not imply an equally large lead on individual airport pairs. Note that national shares use passenger counts, while airport-pair shares use scheduled-flight frequencies, reflecting the data available for each.

\textbf{Rivals' capacity buffer.} Table~\ref{tab:o4-capacity-buffer} shows that rivals' unused capacity relative to IndiGo's traffic fell from 9.79\% in December 2024 to 8.19\% in December 2025. This provides a rough aggregate measure of the spare capacity elsewhere in the market relative to IndiGo's traffic, indicating a smaller buffer in December 2025. However, unused seats would also need to match affected passengers' routes and travel dates, so the ratio cannot directly estimate how many passengers could have been accommodated.
\end{mergedcontribution}

\FloatBarrier
\clearpage
\subsection{Hypotheses revisited}
\label{sec:review-hypothesis-outcomes}
Table~\ref{tab:hypotheses} links the hypothesis outcomes to the results in Section~5. Full hypothesis statements are given in Section~\ref{sec:hypotheses}. H1--H2 concern ownership effects, H3--H4 concern exits, and H5--H7 concern concentration and capacity. O3's entrant-scale expectation is discussed in Section~\ref{sec:review-o3-interpretation}.
\begin{resultstable}{Hypothesis outcomes and supporting evidence.}{tab:hypotheses}
\centering
\fontsize{10}{12}\selectfont\setstretch{1}\setlength{\tabcolsep}{4pt}\renewcommand{\arraystretch}{1.15}
\begin{tabular}{|>{\raggedright\arraybackslash}p{\dimexpr0.23\linewidth-2\tabcolsep-1.3333333333333333\arrayrulewidth\relax}|>{\raggedright\arraybackslash}p{\dimexpr0.23\linewidth-2\tabcolsep-1.3333333333333333\arrayrulewidth\relax}|>{\raggedright\arraybackslash}p{\dimexpr0.54\linewidth-2\tabcolsep-1.3333333333333333\arrayrulewidth\relax}|}
\hline
\textbf{Hypothesis (objective)} & \textbf{Outcome} & \textbf{Supporting evidence from Section 5}\\
\hline
\textcolor{authorSowmya}{\textbf{H1 (O1)}\newline Lasting event steps} & \textcolor{authorSowmya}{Supported for exits and acquisition, weaker for legal mergers} & \textcolor{authorSowmya}{Drift changes are $+531$, $+367$ and $+693$ points for Jet, Tata and Go First. The fitted legal-merger step is smaller, at 137 points ($t=2.5$). Tables~\ref{tab:o1-drift} and~\ref{tab:o1-regression}.}\\
\hline
\textcolor{authorSowmya}{\textbf{H2 (O1)}\newline Partial-ownership wedge} & \textcolor{authorSowmya}{Partially supported, small magnitude} & \textcolor{authorSowmya}{The pre-2022 mean wedge is 22.9 points (0.8\% of Ownership-HHI), with a maximum of 76.6 (2.2\%); it vanishes from January 2022. Table~\ref{tab:review-o1-wedge}.}\\
\hline
\textcolor{authorMeena}{\textbf{H3 (O2)}\newline Exit absorption} & \textcolor{authorMeena}{Partially supported} & \textcolor{authorMeena}{IndiGo leads passenger absorption after both exits (37.7\% and 66.0\%). SpiceJet leads capacity absorption after Jet (30.4\%, versus IndiGo\textquotesingle s 28.5\%). Table~\ref{tab:o2-absorption}.}\\
\hline
\textcolor{authorMeena}{\textbf{H4 (O2)}\newline Contribution to the HHI rise} & \textcolor{authorMeena}{Supported} & \textcolor{authorMeena}{IndiGo adds 473 of 444 passenger-based HHI points after Jet and 632 of 690 after Go First. Capacity-based results are 352 of 354 and 599 of 653 points. Tables~\ref{tab:o2-absorption} and~\ref{tab:review-o2-capacity-contribution}.}\\
\hline
\textcolor{authorGoutham}{\textbf{H5 (O4)}\newline National concentration} & \textcolor{authorGoutham}{Supported} & \textcolor{authorGoutham}{All conditions hold in 20/20 focus-period months. Minimum CR2 is 89.13\% and the minimum share gap is 30.01 pp. Table~\ref{tab:o4-national-results}.}\\
\hline
\textcolor{authorGoutham}{\textbf{H6 (O4)}\newline Route concentration} & \textcolor{authorGoutham}{Supported} & \textcolor{authorGoutham}{HHI exceeds 1,800 in 1,143/1,150 airport-pair/month observations (99.39\%), more than half of the available sample. Table~\ref{tab:o4-route-results}.}\\
\hline
\textcolor{authorGoutham}{\textbf{H7 (O4)}\newline Rivals\textquotesingle\ capacity buffer} & \textcolor{authorGoutham}{Supported} & \textcolor{authorGoutham}{The ratio falls from 9.79\% in December 2024 to 8.19\% in December 2025, a decline of 1.60 pp. Table~\ref{tab:o4-capacity-buffer}.}\\
\hline
\end{tabular}
\resultstablenotes{O denotes the research objective; pp means percentage points. H1 uses the segmented-regression evidence and includes a weaker legal-merger result. H2's assessment reflects the small size of the partial-ownership effect. H5--H7 are descriptive checks, not significance tests or causal estimates. H5 uses December 2024--July 2026; H6 uses April 2019--September 2026. The H7 denominator is all IndiGo monthly RPK.}{Authors' results reported in Section~5.}
\end{resultstable}
\clearpage
\subsection{Connection to the literature review}
\begin{mergedcontribution}{authorSowmya}
\subsubsection{O1: Ownership structure and partial stakes}
Measuring concentration by owner requires knowing who controls each airline and who holds minority stakes in rivals \parencite{obriensalop2000}. Only two airlines had such cross-holdings before the Tata Group placed them in the Air India Group in January 2022 (schedule: Table~\ref{tab:o1-stakes}, Appendix~\ref{app:o1-stakes}).

\textbf{Vistara.} Tata Sons held 51\% from the start of the sample \parencite{sia2015annual,sia2022merger}. The other 49\% belonged to a foreign partner with no stake in any other Indian airline, so it adds no cross-holding term; it enters Equation~O1.1b only through the denominator of $\lambda$, diluting Tata's weight in Vistara's decisions.

\textbf{AirAsia India.} Tata Sons started at 30\% beside AirAsia Investment (49\%) and Telestra (21\%) \parencite{airasia2013jv}, rose to 41.06\% in 2015 and 49\% in March 2016 \parencite{tata2016airasia,airasia2018q4}, held 51\% by 2020, and reached 83.67\% after the December 2020 sale \parencite{airasia2021disposal,airasia2021results}.

\textbf{Inference.} Tata is the only owner with stakes in two rivals, so it alone creates the MHHI wedge: the more of AirAsia India it owns, the more Vistara internalises that rival's profit. The wedge therefore grows with Tata's stake, from about 1 point in 2015 (30\%) to 58 in 2021 (83.67\%), peaks at 76.6 points, and is zero from January 2022 when Tata controls both. At 0.8\% of Ownership-HHI on average, partial ownership matters far less than the full-control grouping of brands. Because stakes are company-reported rather than observed in operating decisions, MHHI measures structure, not conduct.
\end{mergedcontribution}

\begin{mergedcontribution}{authorMeena}
\subsubsection{O2: Concentration after airline exits}
The exit results match \textcite{grosche2020concentration}, who find that concentration in German air transport rose after Air Berlin's bankruptcy and that the largest group, Lufthansa, became more dominant, as IndiGo did here.
\end{mergedcontribution}

\begin{mergedcontribution}{authorSweta}
\subsubsection{O3: Entry, scale and viability}
These findings echo Oliveira and Oliveira (2018), who argue that market growth can attract new entrants but does not, by itself, ensure effective competition. Wang et al. (2018) associate low-cost competition in India with lower fares and higher demand, which is consistent with Akasa Air's rapid growth as a low-cost carrier. However, the present analysis focuses on entrant outcomes rather than estimating the causal effect of entry on fares or demand. Jayathilakan et al. (2024) find a positive association between revenue and passenger volume, but the cases examined here show that revenue growth does not necessarily translate into improved margins or long-term survival.

The findings also highlight the importance of how competition is measured. Brand-level measures, such as those discussed by Kumari et al. (2026), may count each entrant as a separate firm without reflecting substantial differences in their actual market presence. In this study, seven of the nine entrants never exceeded 1\% of national passenger traffic. Ownership and group affiliation also matter when interpreting concentration: Vistara was a Tata-linked joint venture in which Singapore Airlines held a 49\% stake (Tata Sons, 2024), and O1 groups it with the Air India Group from January 2022. Treating brands as independent competitors may therefore produce a different picture from one based on ownership groups. This distinction matters because brand-level market shares do not necessarily represent independent ownership-level competition, particularly after consolidation.
\end{mergedcontribution}

\begin{mergedcontribution}{authorGoutham}
\subsubsection{O4: Concentration and resilience}
\textcite{kumari2026soaring} examine how market concentration and operating efficiency relate to airline performance. Our study complements this work by comparing concentration nationally and on the busiest airport pairs. \textcite{gupta2025antitrust} discuss the risks of dependence on a dominant airline. The capacity-buffer comparison in Table~\ref{tab:o4-capacity-buffer} adds a quantitative measure of rivals' aggregate spare capacity relevant to that concern.
\end{mergedcontribution}

